\section{TRUNG VỊ VÀ TỨ PHÂN VỊ CỦA MẪU SỐ LIỆU GHÉP NHÓM}

\subsection{LÝ THUYẾT CẦN NHỚ}
Cho mẫu số liệu ghép nhóm
\begin{center}
	\begin{tabular}{|c|c|c|c|c|c|}
		\hline
		Nhóm	&$\left[u_1;u_2 \right)$ & $\ldots$ & $\left[u_m;u_{m+1} \right)$ & $\ldots$ &$\left[u_k;u_{k+1} \right)$ \\
		\hline
		Tần số 	& $n_1$ & $\ldots$ & $n_m$ & $\ldots$ & $n_k$ \\
		\hline
	\end{tabular}	
\end{center}

\subsubsection{Trung vị $M_e$ của mẫu số liệu ghép nhóm}
\begin{enumerate}[\iconMT]
	\item \indam{Công thức tính:} 
	\begin{listEX}[1]
		\item [\ding{173}] Giả sử nhóm $ \left[ u_m ; u_{m+1}\right) $ chứa trung vị;
		\item [\ding{174}] $ n_m $ là tần số của nhóm chứa trung vị;
	\end{listEX}
		Khi đó \boxmini{$ M_e = u_m+\dfrac{\dfrac{n}{2}-(n_1+n_2+ \cdots +n_{m-1})}{n_m}\cdot \left( u_{m+1}-u_m\right)$}

		\item \indam{Ý nghĩa:} Trung vị của mẫu số liệu ghép nhóm xấp xỉ cho trung vị của mẫu số liệu gốc, nó chia mẫu số liệu thành hai phần, mỗi phần chứa $50\%$ giá trị.
\end{enumerate}

\subsubsection{Tứ phân vị của mẫu số liệu ghép nhóm}
\begin{itemize}
	\item Gọi $Q_r$ là tứ phân vị thứ $r$ của mẫu số liệu ghép nhóm ($r=1,2,3$);
	\item Nhóm $ \left[ u_m ; u_{m+1}\right) $  chứa tứ phân vị thứ $r$ và $ n_m $ là tần số của nhóm chứa tứ phân vị thứ $r$;
\end{itemize}
Khi đó ta có công thức tính $Q_r$ như sau:
\boxmini{$ Q_r=u_m+\dfrac{r\cdot\dfrac{n}{4}-(n_1+n_2+ \cdots +n_{m-1})}{n_m}\cdot \left( u_{m+1}-u_m\right)$}

\begin{note}
	 Các tứ phân vị của mẫu số liệu ghép nhóm xấp xỉ cho các tứ phân vị của mẫu số liệu gốc, chúng chia mẫu số liệu gồm $4$ phần, mỗi phần chứa $25\%$ giá trị.
\end{note}




